% Generated draft: requires mathematical and source review.
\documentclass[11pt]{article}
\usepackage[margin=25mm]{geometry}
\usepackage{fontspec}
\setmainfont{Latin Modern Roman}
\usepackage{amsmath,amssymb,mathtools,unicode-math}
\setmathfont{Latin Modern Math}
\usepackage{enumitem,needspace,xurl,hyperref}
\setlength{\parindent}{0pt}
\setlength{\parskip}{5pt}
\setlength{\emergencystretch}{3em}
\newcommand{\N}{\mathbb{N}}
\newcommand{\Z}{\mathbb{Z}}
\newcommand{\Q}{\mathbb{Q}}
\newcommand{\R}{\mathbb{R}}
\newcommand{\C}{\mathbb{C}}
\newcommand{\F}{\mathbb{F}}
\newcommand{\A}{\mathbb{A}}
\newcommand{\PP}{\mathbb{P}}
\newcommand{\E}{\mathbb{E}}
\newcommand{\eps}{\varepsilon}
\newcommand{\dd}{\,\mathrm{d}}
\begin{document}
\section*{UnsolvedMath ID: 701724}
\textbf{Source problem code:} catalogue entry
\subsection{Definitions and mathematical statement}
For a finite group \(H\), write \(\pi(H)=\{r:r\text{ is prime and }r\mid |H|\}\). A group is nilpotent if it has a finite central series, and is soluble if its derived series terminates at the trivial subgroup. For subgroups \(A,B\leq G\), let \(AB=\{ab:a\in A,\ b\in B\}\).

The supplied record calls a finite group \(S\) a Schmidt \((p,q)\)-group when \(S\) is not nilpotent, every proper subgroup of \(S\) is nilpotent, and a Sylow \(p\)-subgroup of \(S\) is normal in \(S\). The supplied definition does not explain the role of \(q\), state that \(p\ne q\), or require \(\pi(S)=\{p,q\}\). These are unresolved qualifications of the terminology; no additional restrictions are imposed silently here.

Using the supplied Schmidt \((p,q)\) terminology, define the directed \(N\)-critical graph of a finite group \(H\) to have vertex set \(\pi(H)\) and edge relation
\[
(p,q)\in E_{Nc}(H)
\quad\Longleftrightarrow\quad
\text{there exists }S\leq H\text{ that is a Schmidt }(p,q)\text{-group}.
\]
The intended admissible ordered prime pairs, including whether diagonal pairs are excluded, require confirmation because of the incomplete specification above.

\textbf{Question.} For every finite group \(G\) and all subgroups \(A,B,C\leq G\) satisfying
\[
G=AB=AC=BC,
\]
is it true that
\[
\Gamma_{Nc}(G)=\Gamma_{Nc}(A)\cup\Gamma_{Nc}(B)\cup\Gamma_{Nc}(C)?
\]
Here graph union means union of vertex sets and union of edge sets, with the vertex sets of the subgroup graphs viewed as subsets of \(\pi(G)\). Equivalently, the requested edge equality is
\[
E_{Nc}(G)=E_{Nc}(A)\cup E_{Nc}(B)\cup E_{Nc}(C).
\]
Indeed, the product hypothesis already gives \(\pi(G)=\pi(A)\cup\pi(B)\), so the vertex sets agree. The Schmidt terminology must be resolved consistently in all four graphs before the intended formulation is fully determined.

The supplied source additionally reports that the equality holds if \(A,B,C\) are soluble. This is retained as a source-reported qualification, without independent verification.
\subsection{Short English statement}
Given a finite group that is the product of each pair among three subgroups, must every directed prime pair witnessed by a Schmidt subgroup of the whole group also be witnessed inside at least one of the three subgroups? The source reports an affirmative answer when all three subgroups are soluble. Its definition of a Schmidt \((p,q)\)-group leaves the role of \(q\) unresolved.
\subsection{Sources}
\begin{itemize}
\item Supplied local UnsolvedMath registry record, permanent ID \(701724\), exact source code \texttt{catalogue entry}, title \emph{Kourovka Notebook Problem 20.66}; upstream revision \texttt{d371ee917ac9647fbb06cd6fc529deec6f3dda82}. The record carries the status label \texttt{open}; this is supplied registry metadata, not an independently verified assessment.
\item \emph{The Kourovka Notebook (Unsolved Problems in Group Theory)}, problem 20.66, attributed in the record to the 20th issue (2022), with the 21st issue edition (2026) identified as \texttt{arXiv:1401.0300v48}. The supplied collection description names E. I. Khukhro and V. D. Mazurov as editors. This bibliographic attribution is supplied evidence; the cited text was not independently consulted.
\item Proposed by V. I. Murashka and A. F. Vasil'ev; proposed year supplied as 2022. The statement that the equality holds for soluble \(A,B,C\) is included in the supplied problem text, without a separate proof reference.
\end{itemize}
\end{document}
